\documentclass[10pt]{article}
\usepackage[a4paper,margin=0.65cm,landscape]{geometry}
\usepackage{multicol}
\setlength{\columnseprule}{1pt}
\usepackage[utf8]{inputenc}
\usepackage{amsmath}
\usepackage{amsfonts}
\usepackage{amssymb}
\usepackage{amsthm}
\usepackage{mathtools}
\usepackage{bm}

\newcommand{\C}{\mathbb{C}} % Complex numbers
\newcommand{\R}{\mathbb{R}} % Real numbers
\newcommand{\N}{\mathbb{N}} % Natural numbers
\newcommand{\Q}{\mathbb{Q}} % Rational Numbers
\newcommand{\Z}{\mathbb{Z}} % Integers
\newcommand{\e}{\mathrm{e}} % Natural Constant
\newcommand{\Id}{\bm{I}} % Identity
\newcommand{\zero}{\bm{0}} % Zero vector
\newcommand{\powerset}[1]{\mathcal{P}\left(#1\right)} % Power set
\newcommand{\pbrac}[1]{\left(#1\right)} % Parenthesis
\newcommand{\sbrac}[1]{\left[#1\right]} % Square brackets
\newcommand{\cbrac}[1]{\left\{#1\right\}} % Curly brackets
\newcommand{\abs}[1]{\left\lvert#1\right\rvert} % Modulus
\newcommand{\norm}[1]{\left\lVert#1\right\rVert} % Norm
\newcommand{\func}[3]{#1: #2 \to #3} % Function
\newcommand{\map}[3]{#1: #2 \mapsto #3} % Mapping
\newcommand{\image}[2]{#2\left[#1\right]} % Image
\newcommand{\preimage}[2]{#2\left[#1\right]^{-1}} % Pre-image
\newcommand{\eval}[3]{\left. #1\right\rvert_{#2 = #3}} % Evaluated at
\newcommand{\dif}[3][]{\frac{\mathrm{d}^{#1}#2}{\mathrm{d}#3^{#1}}} % Derivative (dy/dx)
\newcommand{\diff}[3][]{\frac{\mathrm{d}^{#1}}{\mathrm{d}#3^{#1}}#2} % Alternate derivative (d/dx f(x))
\newcommand{\pardif}[3][]{\frac{\partial^{#1} #2}{\partial #3^{#1}}} % Partial derivative
\newcommand{\pardiff}[3][]{\frac{\partial^{#1}}{\partial #3^{#1}}#2} % Alternate partial derivative
\newcommand{\indefint}[2]{\int #1\,\mathrm{d}#2} % Indefinite integral
\newcommand{\defint}[4]{\int_{#1}^{#2}#3\,\mathrm{d}#4} % Definite integral
\newcommand{\tspose}[1]{#1^{\mathrm{T}}} % Transpose
\newcommand{\refto}{\xrightarrow[\mathrm{Elimination}]{\mathrm{Gaussian}}} % Row-echelon form
\newcommand{\rrefto}{\xrightarrow[\mathrm{Elimination}]{\mathrm{Gauss-Jordan}}} % Reduced row-echelon form
\newcommand{\ero}[2][]{\xrightarrow[#1]{#2}} % Elementary row operation
\newcommand{\im}{\mathrm{i}}

\begin{document}
    \begin{multicols*}{3}
        \begin{itemize}
            \item Let $\Z_n \coloneqq \left\{0, 2, \cdots, n - 1\right\}$ and define 
            \begin{equation*}
                U_n \coloneqq \left\{\exp\left(\frac{2k\pi\im}{n}\right) \colon k \in \Z_n\right\},
            \end{equation*}
            then $\left(U_n, \cdot\right)$ is the \textbf{$n$-th root of unity group}.
            \item \textbf{Cancellation laws}: $\forall a, b, u, v \in G$, if $au = av$, then $u = v$; if $ub = vb$, then $u = v$.
            \item $\mathrm{GL}_n\left(\mathbb{F}\right) \coloneqq \left\{\bm{A} \in \mathcal{M}_{n \times n} \colon \det\left(\bm{A}\right) \neq 0\right\}$.
            \item $\forall$ finite $\Sigma$ with $\abs{\Sigma} = n$, define $S_{n, \Sigma}$ to be all bijections from $\Z_{n}^+$ to $\Sigma$, then $\abs{S_{n, \Sigma}} = n!$. (Induction with disjoint union on image of $k + 1$.) $\implies \abs{S_n} = n!$.
            \item $S_n$ is non-Abelian iff $n \geq 3$ ($\implies$: $\exists ab \neq ba$ but neither is identity; $\impliedby$: $\left(12\right)$ and $\left(13\right)$ don't commute)
            \item \textbf{Dihedral}: $r^n = s^2 = 1$, $rs = sr^{-1}$.
            \item $S_n$ is generated by transpositions (induction, define $\tau_{\sigma}$ to swap $k + 1$ and $\sigma\left(k + 1\right)$ if they are unequal and $\mathrm{id}$ otherwise. $\tau_{\sigma} \circ \sigma \in S_k$ and $\tau_{\sigma} = \tau_{\sigma}^{-1}$.)
            \item $S_n$ is generated by $\left(1, k\right)$'s (every cycle is $\prod_{i = 2}^m\left(1, i\right)$, so cycle decomp then transposition decomp).
            \item $S_n = \langle\left(i, i + 1\right)\textrm{'s}\rangle$ ($\left(1n\right) = \left(n - 1, n\right)\left(1, n - 1\right)\left(n, n - 1\right)$)
            \item $S_n = \langle\left\{\left(12\right), \left(12\cdots n\right)\right\}$
            \item \textbf{reversal}: $a < b$ but $\sigma\left(a\right) > \sigma\left(b\right)$. $R\left(\sigma\right)$ is all reversals.
            \item An $m$-cycle has exactly $\left(m - 1\right)$ reversals.
        \end{itemize}
        \textbf{Homomorphism}
        \begin{itemize}
            \item \textbf{Sign}: $\epsilon\left(\sigma\right) = \left(-1\right)^{\abs{R\left(\sigma\right)}}$. even and odd permutations.
            \item Injective iff $\ker = 1$, surjective iff $\mathrm{im} = H$.
            \item \textbf{Alternating group}: $A_n \coloneqq \ker{\epsilon}$. $A_1 = S_1$, $\abs{A_n} = \frac{n!}{2}$.
            \item Let $\theta \colon \Delta \to \Omega$ be a bijection, then $\varphi_{\theta} \colon S_{\Delta} \to S_{\Omega}$ defined by $\varphi_{\theta}\left(\sigma\right) = \theta \circ \sigma \circ \theta^{-1}$ is an isomorphism ($\varphi_{\theta^{-1}} = \varphi_{\theta}^{-1}$) $\implies \varphi_{\theta \circ \eta} = \varphi_{\theta} \circ \varphi_{\eta}$, $\varphi_{\mathrm{id}_{\Delta}} = \mathrm{id}_{S_{\Delta}}$.
        \end{itemize}
        \textbf{Subgroups}
        \begin{itemize}
            \item \textbf{Universal property}: if $H \leq G$, $\forall$ homomorphism $\varphi \colon G \to H$, $\exists!$ homomorphism $\psi \colon G \to K$ where $\mathrm{im}\,{\varphi} \leq K \leq H$ such that $\varphi = i_{K} \circ \psi$.
            \item \textbf{Subgroup criterion}: if $H \subseteq G$, then $H \leq G$ iff $H \neq \varnothing$ and $xy^{-1} \in H$ for all~$x, y \in H$.
            \item Every subgroup of $\Z$ is $n\Z$ ($n$ is min $+$int in subgroup).
            \item \textbf{Finite subgroup criterion}: if finite $H \subseteq G$, then $H \leq G$ iff $H \neq \varnothing$ and $xy \in H$ for all~$x, y \in H$.
            \item Subgroup is closed under intersection.
            \item \textbf{Generated subgroup}: $\left\langle A \right\rangle \coloneqq \bigcap_{A \subseteq H \leq G}H$
            \item $\left\langle A \right\rangle$ is all products of $A$'s elems (use subgrp criterion).
        \end{itemize}
        \textbf{Cyclic groups}
        \begin{itemize}
            \item \textbf{Residue class}: $\bar{a} \coloneqq a + n\Z = \left\{a + kn \colon k \in \Z\right\}$
            \item \textbf{Reduction-modulo-$n$}: $\pi\left(a\right) = \bar{a}$ uniq s.t. $\pi\left(1\right) = \bar{1}$
            \item \textbf{Universal property of $\Z$}: $\forall G$ and $\forall x \in G$, $\exists!$ homomorphism $\varphi \colon \Z \to G$ such that $\varphi\left(1\right) = x$.  For $\Z \slash n\Z$ it's analogous.
            \item $\left(x^a\right)^b = x^{ab}$ (let $\psi\left(1\right) = x$ and $\varphi_k\left(n\right) = kn$, then $\left(\psi \circ \varphi_k\right)\left(1\right) = x^k$. $\left(x^a\right)^b = \left(\psi \circ \varphi_a\right)\left(b\right) = \psi\left(1\right)^{ab}$)
            \item $\Z \slash n\Z = \langle\bar{1}\rangle$ is of order $n$.
            \item Fix $x \in G$ and let $\varphi \colon \Z \to G$ be the unique homomorphism with $\varphi\left(1\right) = x$, then $\mathrm{im}\,{\varphi} = \left\langle x \right\rangle \leq G$ and 
            \begin{equation*}
                \ker{\varphi} = \begin{cases}
                    n\Z & \quad\textrm{if } \abs{x} = n \in \Z^+ \\
                    \left\{0\right\} & \quad\textrm{otherwise}
                \end{cases}.
            \end{equation*}
            ($\ker{\varphi} \subseteq n\Z$ by division algo)
            \item $\abs{x} = \infty \implies\abs{\left\langle x \right\rangle} = \infty$; $\abs{x} = n \implies \abs{\left\langle x \right\rangle} = n$. (Let $\varphi \colon \Z \to G$ be the homomorphism with $\varphi\left(1\right) = x$. Last $x^N$ so $x^{N + 1} = x^k$ for smaller $k$. $\varphi\left(N + 1 - k\right) = 1 \implies N + 1 = k$; $\abs{x} = n \implies x^{nk} = \varphi\left(nk\right) = 1 \implies$ $x^r$'s are pair-wise distinct)
            \item All infinite cyclic groups $\cong \Z$
            \item Every subgroup of an infinite cyclic group is $\langle x^n \rangle$ for $n \in \N$ ($\varphi$ unique isomorphism with $\varphi\left(1\right) = x$, $\psi\left(K\right) = \varphi\left[K\right]$ is bijective, every subgroup of $\Z$ is $n\Z$).
            \item $\forall$ group, if $x^a = x^b = 1$, then $x^{\gcd\left(a, b\right)} = 1$ (by Bezout). 
            \item $\forall$ group and $x \in G$, $\exists n \in \Z$ s.t. $x^n = 1$ iff $\abs{x} \mid n$.
            \item $\forall$ group $G$ with $x \in G$ s.t. $x^n = 1$ for some $n \in \Z$. Let $\varphi \colon \Z \slash n\Z$ be the unique homomorphism s.t. $\varphi\left(\bar{1}\right) = x$, then $\mathrm{im}\,{\varphi} = \left\langle x \right\rangle \leq G$, $\ker{\varphi} = \left\langle\abs{x}\right\rangle \leq \Z \slash n\Z$ 
            \item Every finite cyclic group $\cong \Z \slash n\Z$.
            \item Every subgrp of an order-$n$ cyclic group is $\langle x^a \rangle$ for $a \mid n$
            \item Every subgroup of a cyclic group is cyclic.
            \item Let $G \coloneqq \left\langle x \right\rangle$ be a cyclic group.
            \begin{enumerate}
                \item If $G$ is infinite, then $G = \left\langle x^a \right\rangle$ iff $a = \pm 1$.
                \item If $G$ is of order $n$, then $G = \left\langle x^a \right\rangle$ iff $\gcd\left(a, n\right) = 1$.
            \end{enumerate}
        \end{itemize}
        \textbf{Quotient groups}
        \begin{itemize}
            \item $g_1H = g_2H \iff g_i \in g_jH \iff g_2^{-1}g_1, g_1^{-1}g_2 \in H$.
            \item $\abs{H \backslash G} = \abs{G \slash H}$ ($f \colon H \backslash G \to G \slash H$ by $f\left(X\right) = X^{-1} \coloneqq \left\{x^{-1} \in G \colon x \in X\right\}$, $g \colon H \slash G \to H \backslash G$ by $g\left(X\right) = X^{-1}$ is inverse)
            \item $H \approx gH \approx Hg$
            \item \textbf{Lagrange's Theorem}: $\abs{G} = \abs{G \colon H}\abs{H}$.
            \item $\langle x \rangle \leq G \implies \abs{x} \mid \abs{G}$ for finite $G \implies x^{\abs{G}} = 1$.
            \item A prime-order group is cyclic.
            \item $\left(\Z \slash n\Z\right)^{\times} = \left\{\bar{a} \in \Z \slash n\Z \colon a \leq n \textrm{ and } \gcd\left(a, n\right) = 1\right\}$ ($\supseteq$: Bezout $\implies \overline{sa} = \overline{sa + tn} = \overline{1}$. $\subseteq$: $x = \bar{a}, y = \bar{b}$ s.t. $xy = 1 \implies n \mid ab - 1 \implies ab - kn = 1$)
            \item For any group $G$, 
            \begin{enumerate}
                \item $C_G\left(a\right) = N_G\left(\left\{a\right\}\right)$ for all $a \in G$;
                \item $C_G\left(A\right) = \bigcap_{a \in A}C_G\left(a\right) \subseteq N_G\left(A\right)$;
                \item $Z\left(G\right) = C_G\left(G\right) = \bigcap_{a \in G}C_G\left(a\right)$;
                \item $N_G\left(G\right) = G$;
                \item $C_G\left(1\right) = G$.
            \end{enumerate}
            \item For every group $G$, the followings are equivalent:
            \begin{enumerate}
                \item $G$ is Abelian;
                \item $Z\left(G\right) = G$;
                \item $C_G\left(a\right) = G$ for all $a \in G$;
                \item $C_G\left(A\right) = N_G\left(A\right) = G$ for all $A \subseteq G$.
            \end{enumerate}
            \item $C_G\left(A\right), N_G\left(A\right) \leq G$.
            \item \textbf{Normal subgroup}: $N \trianglelefteq G$ iff $gNg^{-1} \subseteq N \forall g \in G$.
            \item \textbf{Quotient group}: $\left(G \slash N, \cdot\right)$ where $N \trianglelefteq G$. Identity is $N$.
            \item \textbf{Universal property}: Let $N \trianglelefteq G$. For any group $H$ and any homomorphism $\varphi \colon G \to H$ with $N \leq \ker{\varphi}$, $\exists! \bar{\varphi} \colon G \slash N \to H$ such that $\varphi = \bar{\varphi} \circ \pi$ ($\bar{\varphi}\left(\xi\right) = \varphi\left(g\right)$ iff $\pi\left(g\right) = \xi$).
            \item \textbf{1st isomorphism thm}: For any homomorphism $\varphi \colon G \to H$, $\exists! \tilde{\varphi} \colon G \slash \ker{\varphi} \to \mathrm{im}\,{\varphi}$ s.t. $\varphi = i_{\mathrm{im}\,{\varphi}} \circ \tilde{\varphi} \circ \pi$ is the canonical factorisation of $\varphi$ (prove injection by $\ker = 1$).
            \item Let $\varphi \colon G \to H$ be any homomorphism, then 
            \begin{itemize}
                \item $\abs{\mathrm{im}\,{\varphi}} = \abs{G \colon \ker{\varphi}}$;
                \item if $G$ is finite, then $\abs{\mathrm{im}\,{\varphi}} \bigm\vert \abs{G}$;
                \item if $G$ and $H$ are finite, then $\abs{\mathrm{im}\,{\varphi}} \bigm\vert \abs{H}$.
            \end{itemize}
            \item \textbf{2nd isomorphism thm}: if $H, K \leq G$ and $K \trianglelefteq G$, then $H \slash \left(H \cap K\right) \cong HK \slash K$. ($K \trianglelefteq HK$. Define $\varphi = \pi_{HK} \circ i_H \colon H \to HK \slash K$. Use 1st isomorphism thm to get isomorphism $\tilde{\varphi}$. First prove it's surjective and $\mathrm{im}\,\varphi = HK \slash K$. Then $\ker{\varphi} = H\cap K$ because $hk \notin K$ for $h \in K - H$.)
            \item \textbf{3rd isomorphism thm}: if $N, K \trianglelefteq G$ s.t. $N \leq K$, then $K \slash N \trianglelefteq G \slash N$ and $G \slash K \cong \left(G \slash N\right) \slash \left(K \slash N\right)$. (Note: $N \trianglelefteq K$ and $K \slash N \leq G \slash N$ for any $gN \in G \slash N$ and $k \in K$, $gkg^{-1} \in K$ and $\left(gN\right)\left(kN\right)\left(g^{-1}N\right) = \left(gkg^{-1}\right)N \in K \slash N$. For each $k$ there is $k'$ s.t. $k = gk'g^{-1}$ so $kN = \left(gN\right)\left(k'N\right)\left(g^{-1}N\right) \implies K \slash N \trianglelefteq G \slash N$. Define $\varphi = \psi \circ \phi$ where $\psi \colon G \to G \slash N$ and $\phi \colon G \slash N \to \left(G \slash N\right) \slash \left(K \slash N\right)$ are quotient homomorphisms, then $\varphi$ is surjective and $\ker{\varphi} = K$. Now use 1st isomorphism thm.)
            \item \textbf{Lattice point isomorphism thm}: Let $N \trianglelefteq G$, $\pi \colon G \to G \slash N$ be the quotient homomorphism, $\mathcal{H}$ be the family of all subgroups of $G$ containing $N$ and $\mathcal{Q}$ be the family of all subgroups of $G \slash N$, then $\varphi \colon \mathcal{H} \to \mathcal{Q}$ defined by $\varphi\left(H\right) \coloneqq \pi\left[H\right]$ is a bijection such that 
            \begin{enumerate}
                \item if $H \trianglelefteq G$, then $\varphi\left(H\right) \trianglelefteq G \slash N$;
                \item if $X \trianglelefteq G \slash N$, then $\varphi^{-1}\left(X\right) \trianglelefteq G$.
            \end{enumerate}
        \end{itemize}
        \textbf{Group actions}
        \begin{itemize}
            \item $\sigma_g = a \mapsto ga$ is bijective and $g \mapsto \sigma_g$ is a homomorphism (prove $\sigma_{g^{-1}} = \sigma_g^{-1}$).
            \item \textbf{Action homomorphism / perm representation}: $\varphi_{\alpha}\left(g\right)\left(a\right) = \alpha\left(g, a\right)$. \textbf{Always injective}
            \item $\forall$ homomorphism $\varphi \colon G \to S_A$, $\alpha_{\varphi} = \left(g, a\right) \mapsto \varphi\left(g\right)\left(a\right)$ is a group action.
            \item $\psi \colon H \to G$ homomorphism, $\varphi \colon G \to S_A$ perm representation, $\alpha_{\psi}\left(h, a\right) = \varphi\left(\psi\left(h\right)\right)\left(a\right)$ is a group action.
            \item \textbf{Cayley's Thm}: every group $\cong$ some subgroup of some symmetric group.
            \item $\alpha \mapsto \varphi_{\alpha}$ is bijective (prove $\varphi \mapsto \alpha_{\varphi}$ is inverse).
            \item if $p$ is min prime factor of $\abs{G}$, $\abs{G \colon H} = p \implies H \trianglelefteq G$. (consider left multiply action $\alpha$ of $G$ on $G \slash H$. Use 1st isomorphism thm: $\exists K \subseteq S_{G \slash H}$ s.t. $G \slash \ker{\varphi_{\alpha}} \cong K$. Use Lagrange: $\abs{G \colon \ker{\varphi_{\alpha}}} \bigm\vert p!$. Prove $\ker{\varphi_{\alpha}} \leq H \leq G$ and $\abs{H \colon \ker{\varphi_{\alpha}}} = \frac{\abs{G \colon \ker{\varphi_{\alpha}}}}{\abs{G \colon H}} \Biggm\vert \frac{p!}{p} = \left(p - 1\right)!$. Use Lagrange again: $\abs{H \colon \ker{\varphi_{\alpha}}} \bigm\vert \abs{G} \implies \abs{H \colon \ker{\varphi_{\alpha}}}$ has no prime factor.)
            \item $\forall \sigma \in S_A$, define $\mathcal{P}\left(\sigma\right) \colon \mathcal{P}\left(A\right) \to \mathcal{P}\left(A\right)$ by $\mathcal{P}\left(\sigma\right)\left(U\right) = \sigma\left[U\right]$, then $f\left(\sigma\right) = \mathcal{P}\left(\sigma\right)$ is a homomorphism.
            \item $\alpha\bigl(\left(\sigma_A, \sigma_B\right), f\bigr) = \sigma_B \circ f \circ \sigma_A^{-1}$ is a group action of $S_A \times S_B$ on $\mathrm{Maps}\left(A, B\right)$
            \item \textbf{Stable}: $A_0 \subseteq A$ is stable under $\alpha$ if $\alpha\left[G \times A_0\right] \subseteq A_0$.
            \item The followings are equivalent:
            \begin{enumerate}
                \item $A_0$ is stable under $\alpha$;
                \item $\varphi_{\alpha}\left[G\right] = \mathrm{Perm}\left(A\right)_{A_0} \coloneqq \left\{\sigma \in S_A \colon \sigma\left[A_0\right] = A_0\right\}$;
                \item $\sigma_g\left(a\right) \in A_0$ for any $g \in G$ and $a \in A_0$.
                \item $\left.\varphi_{\alpha}\right\rvert_{G \times A_0}$ is an action group homomorphism to $A_0$.
            \end{enumerate}
            \item \textbf{Orbit}: $G \cdot a$ is all $b$ s.t. $\exists g$ s.t. $b = ga$.
            \item \textbf{Transitive action}: only one orbit
            \item \textbf{Stabiliser}: $G_a \coloneqq \left\{g \in G \colon ga = a\right\} \subseteq G$
            \item $f \colon G \slash G_a \to G \cdot a$ by $f\left(gG_a\right) = ga$ is a bijection $\implies \abs{G \colon G_a} = \abs{G \cdot a}$
            \item \textbf{Fixed point}: $ga = a \forall g \in G$. $A^G$ is all fixed points.
            \item $\abs{A} = \abs{A^G} + \sum\abs{G \colon G_{a_i}}$ where $G \cdot a_i$'s non-trivial orbits
            \item Conjugate action: fixed points $Z\left(G\right)$, orbit $C_G\left(g_i\right)$
            \item \textbf{$p$-group}: $\abs{G} = p^{\alpha}$ for some $\alpha \geq 1$. 
            \item \textbf{Sylow $p$-subgroup}: maximal $p$-subgroup
            \item Let $\abs{G} = p^{\alpha}m$, $\alpha \geq 1$ and $p \nmid m$, then $H \leq G$ is a Sylow $p$-subgroup iff $\abs{H} = p^{\alpha}$ or $\gcd\left(\abs{G \colon H}, p\right) = 1$.
            \item If $H$ is a $p$-group acting on $A$, then $\abs{A^H} \equiv \abs{A} \mod p$.
            \item If $H$ is a $p$-subgroup of a finite group $G$, then $\abs{N_G\left(H\right) \colon H} \equiv \abs{G \colon H} \mod p$ (use left multiply action of $H$ on $G \slash H$ and prove $\left(G \slash H\right)^H = N_G\left(H\right) \slash H$).
            \item Centre of any $p$-group is non-trivial
            \item \textbf{Cauchy's Thm}: For every finite group $G$ and any prime $p \bigm\vert \abs{G}$, $\exists g \in G$ s.t. $\abs{g} = p$.
            \item \textbf{Sylow's Theorems}:
            \begin{enumerate}
                \item $\forall$ finite group $G$ and prime factor $p$ with multiplicity $\alpha$ of $\abs{G}$, $\exists$ Sylow $p$-subgroup of order $p^{\alpha}$.
                \item $\forall P, Q \in \mathrm{Syl}_p\left(G\right)$, $\exists g \in G$ such that $Q = gPg^{-1}$.
                \item $\forall G$ finite group with $\abs{G} = p^{\alpha}m$ with $\alpha \geq 1$, and $p \nmid m \geq 1$, $n_p\left(G\right) \mid m$; $n_p \equiv 1 \mod p$; $n_p = \abs{G \colon N_G\left(P\right)}$ for any Sylow $p$-subgroup.
            \end{enumerate}
            \item Every $p$-subgroup is in some Sylow $p$-subgroup
            \item \textbf{Simple group}: non-trivial and $1$ and $G$ are only normal subgroups.
            \item Every group of order $< 60$ is either Abelian simple and $\cong \Z \slash p\Z$ of prime order $p$, or not simple.
            \item $p$-group with order $> p$ is not simple (centre is normal)
            \item Any finite group of order $p^{\alpha}q^{\beta}$ where $p$ and $q$ are primes, $\alpha, \beta \geq 1$, and $p^k \not\equiv 1 \mod q$ for any $k = 1, 2, \cdots, \alpha$ is not simple.
        \end{itemize}
    \end{multicols*}
\end{document}